%%%%%%%%%%%%%%%%%%%%%%%%%%%%%%%%%%%%%%%%%%%%%%%%%%%%%%%%%%%%%%%%%%%%%
%% This is a "template" model document for submission to the
%% American Chemical Society (ACS).
%%
%% The guidance here contains information about how you may wish to
%% modify it to match the requirements of various ACS journals. The
%% ACS do *not* typeset accepted articles using LaTeX, so there is
%% no specific class required.
%%
%% This template deliberately does *not* seek to reproduce
%% the layout of the typeset journal: this is explicitly not
%% required by the ACS for LaTeX submissions.
%%
%% Please report any issues with the template at
%% https://github.com/josephwright/acs-template/issues
%%
%% Released under the Creative Commons 0 license
%% https://creativecommons.org/public-domain/cc0/
%% 
%% Copyight (c) 2025 Joseph Wright
%%%%%%%%%%%%%%%%%%%%%%%%%%%%%%%%%%%%%%%%%%%%%%%%%%%%%%%%%%%%%%%%%%%%%
\documentclass[letterpaper]{article}

%%%%%%%%%%%%%%%%%%%%%%%%%%%%%%%%%%%%%%%%%%%%%%%%%%%%%%%%%%%%%%%%%%%%%
%% Font setup - delete if you are using LuaLaTeX
%%%%%%%%%%%%%%%%%%%%%%%%%%%%%%%%%%%%%%%%%%%%%%%%%%%%%%%%%%%%%%%%%%%%%
\usepackage[T1]{fontenc}
\usepackage{ctex}

%%%%%%%%%%%%%%%%%%%%%%%%%%%%%%%%%%%%%%%%%%%%%%%%%%%%%%%%%%%%%%%%%%%%%
%% Adjust the margins and allow for line spacing
%%%%%%%%%%%%%%%%%%%%%%%%%%%%%%%%%%%%%%%%%%%%%%%%%%%%%%%%%%%%%%%%%%%%%
\usepackage{geometry}
\geometry{margin = 1in}
\usepackage{setspace}

%%%%%%%%%%%%%%%%%%%%%%%%%%%%%%%%%%%%%%%%%%%%%%%%%%%%%%%%%%%%%%%%%%%%%
%% Reference support
%%
%% The recommended method for producing the reference section is
%% to use biblatex. If you wish to use a classical BibTeX
%% approach, this is easiest to achieve using the achemso package.
%% In that case, you should remove the biblatex lines.
%%%%%%%%%%%%%%%%%%%%%%%%%%%%%%%%%%%%%%%%%%%%%%%%%%%%%%%%%%%%%%%%%%%%%
% You can adjust the printing of DOI, article title, etc. using
% package options, e.g. "doi = true"; see the biblatex manual for
% details of adjusting the number of authors printed, e.g.
% "maxnames = 15" to print no more than 15 authors.
\usepackage[style = chem-acs]{biblatex}
\addbibresource{references.bib}
% If you are using classical BibTeX, remove the above lines and 
% uncomment:
%\usepackage{achemso}

%%%%%%%%%%%%%%%%%%%%%%%%%%%%%%%%%%%%%%%%%%%%%%%%%%%%%%%%%%%%%%%%%%%%%
%% Graphic inclusion and scheme and chart support
%%%%%%%%%%%%%%%%%%%%%%%%%%%%%%%%%%%%%%%%%%%%%%%%%%%%%%%%%%%%%%%%%%%%%
\usepackage{graphicx}
\usepackage{float}
% 补充表跨页排版及历史稿 Unicode 字符兼容；均使用现有 TeX Live 宏包。
\usepackage{array,longtable,booktabs,newunicodechar}
\newunicodechar{–}{\textendash}
\newunicodechar{—}{\textemdash}
\newunicodechar{‐}{-}
\newunicodechar{ń}{\'{n}}
\newunicodechar{μ}{\ensuremath{\mu}}
\newunicodechar{γ}{\ensuremath{\gamma}}
\newunicodechar{δ}{\ensuremath{\delta}}
\newunicodechar{α}{\ensuremath{\alpha}}

\usepackage{xcolor}
\usepackage{pgfplots}
\pgfplotsset{compat = 1.18}
\definecolor{acsblue}{RGB}{72,116,176}
\definecolor{acsred}{RGB}{196,78,82}
\definecolor{acsorange}{RGB}{221,132,82}
\definecolor{acsgray}{RGB}{88,89,91}
\definecolor{acsplotbg}{RGB}{239,241,247}
\pgfplotsset{
  acs axis/.style={
    axis background/.style={fill=acsplotbg},
    axis line style={draw=acsgray, line width=0.45pt},
    tick style={draw=acsgray, line width=0.4pt},
    tick label style={font=\footnotesize},
    label style={font=\small},
    title style={font=\small\bfseries},
    grid=both,
    major grid style={draw=white, line width=0.55pt},
    minor grid style={draw=white, line width=0.3pt},
    minor tick num=1,
    scaled ticks=false,
    legend style={draw=none, fill=none, font=\footnotesize},
    every axis plot/.append style={line width=0.85pt}
  },
  acs xbar/.style={
    acs axis,
    xbar,
    xmin=0,
    width=0.88\linewidth,
    height=6.1cm,
    bar width=8pt,
    enlarge y limits=0.16,
    nodes near coords,
    nodes near coords align={horizontal},
    nodes near coords style={font=\scriptsize, xshift=1.2pt}
  },
  acs ybar/.style={
    acs axis,
    ybar,
    width=0.92\linewidth,
    height=6.1cm,
    bar width=11pt,
    enlarge x limits=0.08,
    nodes near coords,
    nodes near coords align={vertical},
    nodes near coords style={font=\scriptsize, yshift=1pt}
  }
}
\newfloat{scheme}{htbp}{los}
\floatname{scheme}{Scheme}
\floatname{chart}{Chart}
\newfloat{graph}{htbp}{loh}

%%%%%%%%%%%%%%%%%%%%%%%%%%%%%%%%%%%%%%%%%%%%%%%%%%%%%%%%%%%%%%%%%%%%%
%% Common support packages
%%%%%%%%%%%%%%%%%%%%%%%%%%%%%%%%%%%%%%%%%%%%%%%%%%%%%%%%%%%%%%%%%%%%%
\usepackage{chemformula} % Formulas using \ch{}
% or
\usepackage[version = 4]{mhchem} % Formulas using \ce{}

%%%%%%%%%%%%%%%%%%%%%%%%%%%%%%%%%%%%%%%%%%%%%%%%%%%%%%%%%%%%%%%%%%%%%
%% Many journals require that sections are unnumbered: this 
%% is activated here
%%%%%%%%%%%%%%%%%%%%%%%%%%%%%%%%%%%%%%%%%%%%%%%%%%%%%%%%%%%%%%%%%%%%%
\setcounter{secnumdepth}{-1}

%%%%%%%%%%%%%%%%%%%%%%%%%%%%%%%%%%%%%%%%%%%%%%%%%%%%%%%%%%%%%%%%%%%%%
%% Place any additional macros here.  Please use \newcommand* where
%% possible, and avoid layout-changing macros (which are not used
%% when typesetting).
%%%%%%%%%%%%%%%%%%%%%%%%%%%%%%%%%%%%%%%%%%%%%%%%%%%%%%%%%%%%%%%%%%%%%
\newcommand*\mycommand[1]{\texttt{\emph{#1}}}

%%%%%%%%%%%%%%%%%%%%%%%%%%%%%%%%%%%%%%%%%%%%%%%%%%%%%%%%%%%%%%%%%%%%%
%% Author and title data:
%% the authblk package is currently the simplest way to provide this
%%%%%%%%%%%%%%%%%%%%%%%%%%%%%%%%%%%%%%%%%%%%%%%%%%%%%%%%%%%%%%%%%%%%%
\usepackage{authblk}
\author[1]{Xinge Xia}
\author[1]{Fan Wu}
\affil[1]{College of Cyber Security, Jinan University, Guangzhou, China}
\author[2]{Xianyou Duan}
\affil[2]{College of Chemistry and Materials Science, Jinan University, Guangzhou, China}
\author[1]{Hongyi Chen}
\author[1]{Bingwen Feng*}

\title{纳米硒细胞毒性预测：材料化学增强特征下的机器学习模型比较}
% Use the \date command for email address(s) of corresponding authors
\date{*通讯作者： bingwfeng@gmail.com}

\begin{document}

\maketitle

\begin{abstract}
  纳米硒（SeNPs）因具有良好的生物相容性和多重生物活性而受到广泛关注，但其毒性效应受粒径、浓度、表面修饰和材料化学状态等因素共同影响，传统实验方法难以快速完成系统评估。本文围绕纳米硒细胞毒性预测问题，构建了“基准数据集模型比较 + 专用数据集增强特征建模”的研究框架。首先，在来源于 Masarkar 等人的多类型纳米颗粒公开数据集 \cite{masarkar2024} 上，对随机森林、XGBoost、基于注意力的随机森林（ABRF）、基于叶注意力的随机森林（LARF）、TabTransformer 和蒸馏 TabTransformer 进行了系统比较；随后基于我们整理的 383 条硒纳米颗粒有效记录，进一步开展专用数据集特征统计、数据预处理分析与模型构建。

结果表明，在大样本异质表格数据上，TabTransformer 在回归任务中取得最高 $R^2$（0.768），而树模型在分类任务中表现更稳定，其中 XGBoost 的 AUC 为 0.945。在硒纳米颗粒专用数据集上，数据来源覆盖 43 篇文献、时间跨度为 2008--2022 年，引入材料化学相关增强信息后，树模型的拟合能力较基础特征设置有所改善，其中 XGBoost 在增强特征设置下取得较优结果，三折交叉验证 MAE 为 7.20\%，$R^2$ 为 0.957；随机森林在增强特征下也达到较高精度（MAE 7.63\%，$R^2$ 0.915）。进一步的自动筛选表明，IC50、元素 Se 含量、长期稳定性和Se 3d 结合能 构成表现较优的增强变量组合；但考虑到 IC50 与细胞毒性终点之间存在天然关联，其贡献更适合视为辅助验证信息，而非与先验理化特征完全同层级的独立预测变量。

本研究说明，面向单一材料体系的毒性预测，仅依赖通用理化参数是不够的，能够刻画材料本征化学属性的增强特征具有重要补充价值。所构建的数据整理、特征统计和建模流程可为纳米硒毒性预测平台开发、实验优先级筛选及后续外部验证研究提供方法基础，而含 IC50 的结果仍需结合真实预测场景进一步审慎解释。

\end{abstract}

\section*{Keywords}

Some journals require keywords: these normally should be given immediately
after the abstract.

\section*{Abbreviations}

Some journals require a list of abbreviations: these normally should be given
immediately after the keywords (if required).

%%%%%%%%%%%%%%%%%%%%%%%%%%%%%%%%%%%%%%%%%%%%%%%%%%%%%%%%%%%%%%%%%%%%%
%% Start the main part of the manuscript here.
%%%%%%%%%%%%%%%%%%%%%%%%%%%%%%%%%%%%%%%%%%%%%%%%%%%%%%%%%%%%%%%%%%%%%
\section{引言}

硒是哺乳动物必需的微量元素，参与谷胱甘肽过氧化物酶等硒蛋白的合成，在维持机体氧化还原稳态中发挥重要作用。与亚硒酸盐、硒酸盐等无机硒相比，零价态纳米硒具有更高的生物利用度和更低的急性毒性。得益于其良好的生物相容性及多种生物学活性，纳米硒在抗肿瘤、抗菌、抗糖尿病及抗氧化等生物医学领域展具有广阔的应用前景\cite{hosnedlova2018nanoselenium,bisht2022senps,sowmya2024nanoselenium,karthik2024senps}。然而，硒的治疗指数（therapeutic index）较低，其有效剂量与毒性剂量之间界限有限，即使在药理活性剂量范围内，纳米硒仍可能产生氧化损伤及细胞毒性等不良反应\cite{wang2020prooxidation}。过量暴露还可能导致发热、恶心，以及呼吸系统、心血管系统、消化系统损伤甚至肾衰竭等系统性毒性反应\cite{hosnedlova2018nanoselenium}，显示其生物安全性评价的重要性与必要性。
 
纳米材料的生物毒性受到粒径、形貌、表面性质及化学组成等多种因素的综合影响，各理化参数之间存在显著的非线性耦合特征，增加了毒性效应预测与机制解析的复杂性\cite{ji2025rosml}。传统的纳米毒性评价主要依赖体外细胞实验以及体内动物实验（如大鼠、小鼠和斑马鱼模型），但该类方法普遍存在实验周期长、成本高及检测通量有限等问题\cite{ahmadi2024toxicity}。随着计算毒理学的发展，纳米定量构效关系（nano-QSAR/QSTR）等方法逐渐成为纳米毒性研究的重要补充工具\cite{roy2023tio2}。近年来，随机森林（RF）、支持向量机（SVM）及神经网络（NN）等机器学习算法已被广泛应用于金属及金属氧化物纳米颗粒的细胞毒性预测\cite{zhang2023qsar}，并结合聚类分析、特征筛选等方法进一步提升模型的稳定性、预测精度及泛化能力\cite{sang2022mlcytotoxicity}。
 
针对纳米硒的研究表明，其毒性效应与粒径、暴露浓度及表面修饰等因素密切相关。粒径方面，随着纳米硒颗粒粒径减小及比表面积增加，其生物活性和细胞摄取效率通常进一步增强，同时潜在毒性风险也随之升高\cite{zhang2023senpsbio}。浓度方面，Wang 等\cite{wang2020prooxidation}的研究表明，大鼠肝细胞暴露于 24~$\mu$M 和 48~$\mu$M SeNPs 24~h 后，细胞存活率分别下降至 87.76\% 和 51.37\%，呈现明显的剂量依赖性毒性效应。表面修饰同样能够显著影响纳米硒的生物安全性与毒性行为。例如，螺旋藻多糖（SPS）、江蓠多糖（GLP）及黄芪多糖（AASP）等表面修饰可增强纳米硒对肿瘤细胞的选择性杀伤作用，并降低其对正常细胞的毒性\cite{zhang2023senpsbio,ji2023astragalus}；而季铵化壳聚糖--没食子酸修饰的 SeNPs 在 Colo357 和 HaCaT 细胞中则表现出较高的细胞毒性\cite{lunkov2023chitosan}。上述研究表明，纳米硒的毒性效应具有显著的理化性质依赖特征，难以通过传统经验进行简单外推，因此有必要建立面向纳米硒特定物化参数的毒性预测模型。
 
近年来，机器学习与深度学习在高维特征挖掘、复杂非线性关系建模及模式识别等方面展现出显著优势，为计算毒理学研究提供了新的技术路径\cite{hosni2025mlnp}。Shirokii 等通过相关性特征筛选及网格搜索优化超参数，将无机纳米材料细胞毒性的定量预测模型在 10 折交叉验证中的 $Q^2$ 提升至 0.86\cite{shirokii2023quantitative}。此外，LightGBM 等集成学习算法在区分细胞毒性与非细胞毒性纳米颗粒方面表现出优于传统机器学习模型的预测性能\cite{hosni2025mlnp}
，Boruta-RF-ARM 等组合建模框架则进一步识别出浓度、核壳材料及细胞类型等关键毒性影响因素\cite{yildirim2025rheumatoid}。尽管现有研究已在纳米毒性预测领域取得一定进展，但相关模型主要集中于有机小分子或通用纳米材料，对纳米硒毒性效应的针对性研究仍相对有限，其专用预测模型与特征建模方法尚不完善。基于此，本研究围绕纳米硒毒性预测构建了由多模型比较、增强特征建模及自动特征筛选组成的分析框架，以浓度、粒径、表面修饰及材料化学属性等关键变量作为输入，系统评估不同模型在通用纳米颗粒数据集与纳米硒专用数据集中的适用性与预测性能，旨在为纳米硒生物安全性评价、实验优先级筛选及后续智能预测平台开发提供方法学基础。


\section{方法}

\paragraph{数据收集。}
研究共采用两个相互独立的数据集对模型进行训练与验证，以充分评估模型架构的可靠性，实现针对纳米硒的细胞毒性预测。

\begin{enumerate}
  \item \textbf{通用数据集}：
该数据集来源于 Masarkar 等人 \cite{masarkar2024} 开展的纳米颗粒细胞毒性元分析研究。
数据集汇总了已发表文献中关于金属、金属氧化物、聚合物以及碳基纳米颗粒四大类别的细胞毒性实验数据，是当前纳米材料机器学习毒性预测领域具有代表性的公开数据资源之一。
由于本研究的模型改进工作是在其原有框架基础上展开，因此采用该数据集对基线模型进行复现，以保证模型优化过程具备良好的可比性与有效性。
该数据集原始共包含 4111 条实验记录，经清洗后保留 4102 个有效样本用于后续训练测试划分与模型比较。

  \item \textbf{自建数据集}：
针对纳米硒这一特定研究对象，本研究从 43 篇发表于 2008--2022 年的同行评议文献中系统提取实验记录，构建纳米硒细胞毒性预测专项数据集。数据整理采用大模型粗提取与人工精校验证相结合的流程：首先利用提示词工程，引导大模型按照 Masarkar 等人\cite{masarkar2024} 的特征体系提取 43 篇文献中的特征属性；对于正文未明确报告而可从图表获得的数据，进行人工数字化提取；最后由研究人员逐项对照原始文献核验全部记录。来源文献完整清单见补充材料表~\ref{tab:source-publications}，逐文献记录数尚待根据原始 383 条数据核实。整理后共包含 383 条有效记录，每条记录对应 37 个特征参数，覆盖纳米硒理化性质、暴露条件及生物学终点，为后续增强预测模型的训练、验证与评价提供数据基础。
\end{enumerate}

为进一步说明自建数据集的文献来源结构，本研究按题名和研究目的将 43 篇纳米硒相关文献划分为四类：材料构建/表面修饰、机制毒理/安全性、递送/联合治疗、诊疗/免疫及其他应用。该分类仅用于展示文献覆盖范围和主题分布，不改变后续建模中对原始实验记录的提取规则。图~\ref{fig:senp-literature-cloud} 展示文献发表年份与主题分类的点云分布。

\begin{figure}[H]
  \centering
  \begin{tikzpicture}
    \begin{axis}[
      acs axis,
      width=0.82\linewidth,
      height=6.5cm,
      xmin=2007.5,
      xmax=2022.8,
      ymin=0.55,
      ymax=4.45,
      xlabel={发表年份},
      xtick={2008,2010,2012,2014,2016,2018,2020,2022},
      ytick={1,2,3,4},
      yticklabels={材料/表面,机制/安全,递送/联合,诊疗/免疫/其他},
      y tick label style={font=\scriptsize},
      title={\textbf{(A)} 43 篇文献点云}
    ]
      \addplot+[only marks, mark=*, mark size=1.9pt, acsblue] coordinates {
        (2008,0.78) (2011,1.10) (2011.88,0.92) (2012.12,1.24)
        (2012.36,0.80) (2013.16,1.08) (2014,0.92) (2019,1.20)
      };
      \addplot+[only marks, mark=square*, mark size=1.8pt, acsred] coordinates {
        (2011.15,1.80) (2012.86,2.08) (2013.05,1.86)
        (2013.30,2.22) (2018,1.92) (2019.15,2.12)
      };
      \addplot+[only marks, mark=triangle*, mark size=2.1pt, acsorange] coordinates {
        (2012.44,2.76) (2012.68,3.18) (2013.42,2.90) (2014.04,3.24)
        (2014.32,2.78) (2015,3.12) (2015.96,2.86) (2016.22,3.22)
        (2017,2.94) (2018.16,3.16) (2019.30,2.84) (2020,3.12)
        (2021,2.90) (2021.24,3.20)
      };
      \addplot+[only marks, mark=diamond*, mark size=2pt, acsgray] coordinates {
        (2016.34,3.78) (2016.54,4.06) (2016.74,4.30) (2018.30,3.88)
        (2018.50,4.18) (2019.44,3.78) (2019.64,4.08) (2019.84,4.30)
        (2020.14,3.92) (2021.34,4.20) (2021.54,3.84) (2022,4.10)
        (2022.18,3.78) (2022.36,4.30) (2022.54,3.96)
      };
    \end{axis}
  \end{tikzpicture}

  \vspace{0.8em}

  \begin{tikzpicture}
    \begin{axis}[
      acs xbar,
      width=0.62\linewidth,
      height=4.6cm,
      xmax=16,
      xlabel={文献数量},
      symbolic y coords={机制/安全,材料/表面,递送/联合,诊疗/免疫/其他},
      ytick=data,
      y tick label style={font=\scriptsize},
      title={\textbf{(B)} 分类数量}
    ]
      \addplot+[fill=acsblue, draw=none] coordinates {(6,机制/安全) (8,材料/表面) (14,递送/联合) (15,诊疗/免疫/其他)};
    \end{axis}
  \end{tikzpicture}
  \caption{自建纳米硒数据集 43 篇来源文献的发表年份点云与主题分类分布。}
  \label{fig:senp-literature-cloud}
\end{figure}

\paragraph{数据预处理。}
对于通用数据集，我们沿用了 Masarkar 等人 \cite{masarkar2024} 的预处理方案：
首先，我们删除了数据集中包含大量缺失值的特征属性（如多分散指数 PDI 和长径比），以及与预测无关的文献信息（如文章 ID、发表年份和 DOI 编号），将特征数量从初始的 30 个精简至 15 个。
其次，对于某些在物理化学性质（如尺寸、Zeta 电位、表面电荷和 PDI）上存在严重信息缺失的纳米颗粒（例如量子点 QD 和羟基磷灰石 HAP），我们剔除了其对应的所有数据点。

随后，我们将剩余的数据集按照材料类型统一划分为四大类纳米颗粒：金属（M）、金属氧化物（MO）、聚合物（P）和碳基（C）纳米颗粒。
在处理具体的物理化学特征时，为了整合不同测试条件下的数据，我们采用了以下优先级策略：

\begin{enumerate}
  \item \textbf{尺寸（Size）特征}：
优先采用无任何介质条件下的尺寸；若该值缺失但提供了水中和培养基中的尺寸，则采用水中的尺寸；若原数据仅提供了一种条件（无介质、水中或培养基中）下的尺寸，则直接使用该已知值。

  \item \textbf{Zeta 电位特征}：
采用类似策略，优先采用水中的 Zeta 电位；若仅提供了一种条件下的 Zeta 电位，则直接使用该值。
\end{enumerate}

在执行上述整合、审查目标变量缺失并严格过滤异常值后，最终清洗后的数据集保留了 4102 个有效数据点和 15 个特征属性。
最后，为了使数据完全适合模型训练，我们使用“min-max”缩放器（MinMaxScaler）对编译后的数据集进行了归一化处理，并对所有的类别属性进行了标签编码（Label encoding）。

对于自建的纳米硒专项数据集，预处理步骤进行了改进。
本研究自建的纳米硒毒性数据集共包含 383 个样本与 37 个特征。
数据集覆盖了纳米硒研究中常见的实验条件区间，包括不同的暴露剂量、粒径范围、形态类型与表面修饰策略，可反映纳米硒与其理化性质的关系在多个维度上的变化。

自建数据集数据预处理包括以下步骤：

\begin{enumerate}
  \item \textbf{缺失值处理}：由于粒径、Zeta 电位、暴露时间、浓度、细胞存活率诸字段都有若干缺失值，故首先直接剔除含缺失值的记录。此外，“检测方法”字段中若干条目“Characterization”“Western”“Cellular Uptake”属于非毒性检测目的条目，故一并删除。最终共剔除 262 条记录，保留 383 个符合要求的有效样本。
  \textbf{上述剔除仅针对所列关键字段；派生特征（如粒径比）和非关键分类变量（如表面电荷）的缺失值在剔除之后单独处理，不在此列。}

  \item \textbf{多报告值的优先级选择}：参照 Masarkar 等人 \cite{masarkar2024} 提出的方法：同一纳米颗粒在文献中报告多种尺寸时，按无介质中的粒径 > 水相粒径 > 培养基中的粒径的顺序选取，而 Zeta 电位宜取水相测量值，若无该数据则取文献中所给的唯一测量值。

  \item \textbf{单位统一与换算}：把各浓度单位都统一换算为 $\mu$g/mL，凡以 $\mu$M 为单位的浓度都根据硒的分子量（78.96 g/mol）换算为 $\mu$g/mL，同时把暴露时间统一转换为小时（h）。

  \item \textbf{数值特征变换}：由于浓度、时间两变量都呈偏态分布，对二者做对数变换：
        \[
          \text{Log\_Conc} = \ln(1 + \text{Concentration\_μg/mL}), \quad
          \text{Log\_Time} = \ln(1 + \text{Time\_h})
        \]
        同时引入“粒径比”= 水合粒径 / 核心粒径来表征颗粒在水相中的聚集倾向，所有缺失值均按 1.5 填充。

  \item \textbf{分类变量编码}：对表面电荷（正/负/中性）、细胞类型、修饰基团、检测方法、晶型（结晶/无定形）等类别特征，采用标签编码转换为数值，其中表面电荷缺失值填充为“未知”。

  \item \textbf{相关性检查}：对数值特征计算 Pearson 相关系数，未发现 $r > 0.8$ 的高相关度特征对，表明特征冗余度很低。
\end{enumerate}

\paragraph{特征工程与多模态融合。}

在此基础上，模型输入特征被组织为基础特征和增强特征两组。基础特征包含以下几类信息：暴露条件（浓度、时间）、颗粒理化属性（粒径、Zeta 电位、尺寸比、表面电荷）、生物体系属性（细胞类型）以及实验条件（检测方法、修饰基团）。
针对硒纳米颗粒专用建模，在基础特征之外进一步引入 5 项增强变量，即 IC50、元素 Se 含量、长期稳定性、Se 3d 结合能 和晶型编码。其中，元素 Se 含量、长期稳定性、Se 3d 结合能 与晶型编码主要用于表征材料本征化学属性；IC50 则与毒理响应终点具有较强关联，因此本文将其视为辅助变量，用于比较含/不含该变量时模型表现的差异，而不将其简单等同于可在真实预测前稳定获得的先验理化特征。

为利用不同特征模态之间的互补信息，本研究采用多模态特征融合策略，将输入划分为三个模态：（1）基础理化描述符，包括浓度、时间、粒径、Zeta 电位、尺寸比、表面电荷、细胞类型、检测方法和修饰基团；（2）材料化学增强描述符，包括元素 Se 含量、长期稳定性、Se 3d 结合能和晶型编码；（3）可选辅助描述符 IC50。融合分为三个阶段：首先拼接基础特征与材料化学增强特征，构建联合表示；随后，在 IC50 可用时以其作为查询，对联合表示计算交叉注意力权重，自适应突出相关特征，而在 IC50 缺失时旁路注意力机制、直接使用联合表示；最后加入残差连接以保留原始联合表示，并将所得多模态特征矩阵输入预测模型。该设计将材料本征化学信息与通用理化描述符整合，并根据辅助信息是否可用动态调整特征权重。

为评估新增增强信息对模型性能的提升作用，保持基础特征不变，对上述 5 项增强变量的全部非空组合进行穷举搜索，共 31 种组合，并以三折交叉验证 MAE 作为筛选指标，用于确定表现较优的增强变量子集。需要指出的是，若最优组合包含 IC50，则其结果更适合作为辅助分析依据，而不宜直接等同于真实部署情境下的预测能力提升。

\paragraph{模型与评价。}
本研究围绕表格型纳米毒性数据比较了树模型、注意力增强树模型与深度表格模型三类方法。在基准数据集上，纳入随机森林、XGBoost、ABRF、LARF、TabTransformer 和蒸馏 TabTransformer 六种模型，以系统评估不同架构对异质纳米颗粒数据的适应性；在硒纳米颗粒专用数据集上，重点比较基础特征与增强特征条件下各模型的性能差异。随机森林作为核心基线模型，用于评估中小样本场景下的稳健性与可解释性；XGBoost 用于测试梯度提升树在增强特征条件下的高精度建模能力；ABRF 与 LARF 用于考察注意力机制引入后是否能够提升树模型表现；TabTransformer 与蒸馏 TabTransformer 则作为深度表格模型代表，用于与传统树模型进行对照。

本文对模型性能采用回归、分类两类指标进行评价。回归任务中采用决定系数（\(R^2\)）、平均绝对误差（MAE）和均方根误差（RMSE）。\(R^2\) 是对模型所拟合的目标变量方差的解释比例的自然度量，理论上不超过 1（当模型表现差于均值基线时可能为负值），越接近 1 表明拟合效果越好。MAE 是预测值与真实值绝对误差的算术平均，其单位与目标变量一致（\%），因此也是预测典型偏差的直接度量。RMSE 对较大误差更敏感，适用于考察预测的稳定性。

分类任务沿用 Masarkar 等人\cite{masarkar2024} 所采用的 85\% 细胞存活率阈值，将样本划分为“高存活”和“低存活”两类，以保持与基准数据集一致并便于直接比较模型表现，继而用准确率（Accuracy）、F1 分数和 ROC 曲线下面积（AUC）作为主要指标。准确率即正确分类的比例，F1 分数是精确率、召回率的调和平均，对类别不平衡问题具有适应性，而 AUC 衡量模型区分正负类的能力，0.5 为随机猜测，1.0 表示完全区分。

对于大规模多类型纳米颗粒基准数据集，采用 7:3 训练测试划分，以测试集性能作为主要比较依据。回归任务使用 $R^2$、MAE 和 RMSE 进行评价，分类任务使用 Accuracy、F1 和 AUC 进行评价。

对于硒纳米颗粒专用数据集，采用三折交叉验证，以 CV MAE 和全量 $R^2$ 作为主要回归评价指标，并用于比较基础特征与增强特征两种设置下不同模型的表现。增强特征筛选阶段同样采用三折交叉验证 MAE 作为统一指标，以保证不同特征组合之间的结果可比性。

\paragraph{预测模型的架构细节。}
为系统评估不同架构在纳米毒性预测任务上的适用性，我们尝试了多种模型，涵盖传统树模型、梯度提升树、注意力增强树模型以及深度表格模型。
经过初步性能比较，最终选择随机森林、XGBoost、ABRF、LARF、TabTransformer 以及蒸馏 TabTransformer 六种模型作为主要分析对象。
下面将依次阐述各模型的设计思想及其在本任务中的定位。

随机森林是以决策树为基础的典型的集成学习方法，通过自助采样（Bootstrap）构建多个决策树，并基于其结果进行回归或分类。该方法具有较强的抗过拟合能力，能有效处理混合型特征（数值型和类别型），并且对缺失值不敏感，能直接估计特征重要性，因此在中小型表格数据任务中呗广泛用作基线模型。
基于此，本文采用随机森林来考察其在样本量有限的硒纳米颗粒数据集上的稳健性及可解释性。


XGBoost 是一种基于梯度提升树的集成学习框架，在模型构建过程中引入了二阶导数信息以及 L1、L2 正则化项，因此在表格数据任务中具有较高的预测性能。该方法通过串行迭代方式不断拟合残差，通常能够获得较高的模型精度，但同时对超参数设置较为敏感，需要通过合理调参降低过拟合风险。
本文引入 XGBoost 模型，旨在分析梯度提升树方法在加入元素组成、稳定性、结合能等增强信息后，对硒纳米颗粒数据集预测性能的变化，并进一步比较含/不含 IC50 时模型表现的差异。

ABRF 是在标准随机森林基础上加入了一个注意力网络来对各棵决策树的预测做动态加权处理：先训练普通的随机森林，再对每个样本收集所有树的预测值，接着用一个小型神经网络直接、自动地学习各棵树就该样本的权重，最后用各树预测值的加权和作为最终预测结果。因此其本质是为了用注意力机制合理地挖掘不同树、不同样本之间预测可靠性的差异。
ABRF 是一种在标准随机森林基础上引入注意力机制的集成学习方法。该方法首先训练传统随机森林模型，并收集各决策树对样本的预测结果；随后利用小型神经网络学习不同决策树在不同样本上的权重分配；最终通过对各树预测结果进行加权求和得到模型输出。核心思想是在样本层面建模不同决策树预测贡献的差异，以提升模型的表达能力。
本文引入 ABRF 模型旨在分析单层注意力机制对树模型预测性能的影响。


LARF 在 ABRF 的基础上进一步引入了特征级注意力机制。在模型训练过程中，首先学习特征权重向量，对原始特征进行加权处理，以增强重要特征在决策树分裂过程中的作用；随后结合树级注意力机制，对不同决策树的预测结果进行动态加权，从而形成包含特征级与树级的双层注意力结构。该方法能够同时优化特征表示与树模型集成过程。
本文引入 LARF 来分析更复杂的注意力机制对纳米硒数据集预测性能的影响。

TabTransformer 是一种面向表格数据的深度学习模型，其核心思想是利用 Transformer 的注意力机制对类别特征进行编码，将其表示为嵌入向量，并与数值特征拼接后输入多层感知机进行预测。由于 Transformer 能够建模类别特征之间的高阶交互关系，因此该模型在包含多种类别变量的表格数据任务中具有一定优势。本文将 TabTransformer 作为深度表格模型的代表，用于评估深度学习方法在纳米毒性数据集上的预测性能，并与传统树模型进行比较。

蒸馏 TabTransformer 则基于知识蒸馏方法，以完整的 TabTransformer 作为教师模型，对参数规模较小的学生模型进行训练。该方法通过降低嵌入维度及多层感知机隐藏层规模，实现模型轻量化，并在训练过程中同时利用教师模型输出的软标签与真实标签信息，以兼顾模型精度与计算效率。
故本文引入该模型，旨在分析资源受限场景下轻量化深度表格模型的部署可行性及其性能表现。

由于六种模型在大规模、多类型纳米颗粒基准数据集上已经进行了较为系统的比较，因此本文重点分析不同模型架构对异质纳米颗粒数据的适应能力。
在纳米硒专用数据集上，由于样本量较小（383 条），本文主要选取随机森林、XGBoost 和 TabTransformer 三种代表性模型进行比较，并进一步分析基础特征与增强特征两种特征设置下的性能差异。其中，随机森林与 XGBoost 用于评估传统树模型的表现，TabTransformer 用于考察深度学习模型在小样本表格数据中的适用性。
对于 ABRF、LARF 以及蒸馏 TabTransformer，由于其模型结构相对复杂、训练成本较高，因此仅在大规模基准数据集上进行评估，以兼顾实验复杂度、计算成本与结果分析的可读性。

\paragraph{可视化特征效应与交互。}
特征重要性分析能够确定对模型预测贡献较大的变量，但仅凭重要性排序难以说明特征与细胞存活率之间的关系形态，例如这种关系是近似线性、单调变化还是存在更复杂的非线性响应。参考 Masarkar 等人 \cite{masarkar2024} 的处理思路，本研究先对纳米硒数据集的样本组成、连续型变量和类别型变量进行可视化描述，以明确模型训练前的数据覆盖范围与潜在偏倚；在完成模型训练后，再结合特征重要性进一步考察单个特征效应与双特征交互关系。对于单特征效应，可通过部分依赖图（partial dependence plot, PDP）和个体条件期望图（individual conditional expectation, ICE）观察模型预测值随某一特征变化的整体趋势及样本间异质性；对于两个特征的联合影响，可进一步采用二维 PDP 等高线图展示其组合变化对预测细胞存活率或毒性类别的影响。

在进行特征效应可视化之前，首先需要明确样本组成和变量分布。本研究数据集共包含 383 条有效记录（剔除粒径、Zeta 电位等关键理化性质缺失以及细胞存活率缺失的记录后），来源于 43 篇已发表文献，文献发表年份覆盖 2008--2022 年。清洗后的有效记录主要集中于 2018 年（157 条，占 41.0\%），其次为 2019 年（96 条，占 25.1\%）和 2012 年（44 条，占 11.5\%）。所有纳米硒颗粒的形貌均为球形，物种来源以人类细胞为主（371 条，占 96.9\%）。

图~\ref{fig:senp-year-distribution} 显示了样本来源年份的集中程度，表~\ref{tab:senp_feature_statistics} 汇总了主要连续型特征的描述性统计量，并为后续可视化解释提供了基础。粒径的分布明显右偏（偏度 3.92），绝大多数颗粒集中在 30 nm 附近，极少数样本超过 200 nm，符合纳米硒常规粒径范围。水合粒径右偏程度较轻（偏度 1.69），其均值（105 nm）明显大于核心粒径均值（65 nm），这表明水化层以及可能存在的弱聚集效应。Zeta 电位分布接近对称（偏度 -0.08），范围从 $-60$ mV 到 $+36$ mV，其中负电位占主导，反映出多数表面修饰（如聚糖、PVP）带负电。

\begin{figure}[H]
  \centering
  \begin{minipage}{0.48\linewidth}
    \centering
    \begin{tikzpicture}
      \begin{axis}[
        acs xbar,
        width=\linewidth,
        height=5.1cm,
        xmax=175,
        xlabel={样本量},
        symbolic y coords={2012,其他年份,2019,2018},
        ytick=data,
        y tick label style={font=\scriptsize},
        title={\textbf{(A)} 来源年份频率}
      ]
        \addplot+[fill=acsblue, draw=none] coordinates {(44,2012) (86,其他年份) (96,2019) (157,2018)};
      \end{axis}
    \end{tikzpicture}
  \end{minipage}
  \hfill
  \begin{minipage}{0.48\linewidth}
    \centering
    \begin{tikzpicture}
      \begin{axis}[
        acs axis,
        width=\linewidth,
        height=5.1cm,
        ymin=0,
        ymax=105,
        ylabel={累计占比（\%）},
        symbolic x coords={2018,2019,其他年份,2012},
        xtick=data,
        x tick label style={rotate=35, anchor=east, font=\scriptsize},
        ymajorgrids=true,
        title={\textbf{(B)} 累计贡献}
      ]
        \addplot+[acsred, mark=square*, mark size=1.7pt] coordinates {(2018,41.0) (2019,66.1) (其他年份,88.6) (2012,100.0)};
      \end{axis}
    \end{tikzpicture}
  \end{minipage}
  \caption{纳米硒数据集中文献来源年份的样本量分布与累计贡献。}
  \label{fig:senp-year-distribution}
\end{figure}

浓度分布极端右偏（偏度 8.28），四分位距仅为 $0.52$--$10$ $\mu$g/mL，但最大值高达 $490$ $\mu$g/mL，提示少数实验采用极高剂量，可能对模型产生杠杆效应。暴露时间呈轻微左偏（-0.42），中位数 72 h，说明多数实验采用 24 h 或 72 h 的暴露方案。细胞存活率分布近似均匀（偏度 -0.17），均值与中位数均接近 50\%，意味着数据集正负样本较为平衡，有利于分类模型训练。

为直观展示上述分布，我们进一步基于描述性统计绘制了连续型特征的可用样本量和偏度对照图（图~\ref{fig:senp-continuous-summary}）。该图不对分位数之外的个体观测分布进行外推，而是直接使用表~\ref{tab:senp_feature_statistics} 中的 $n$ 和偏度刻画数据覆盖度与分布偏斜程度。

\begin{table}[H]
  \centering
  \caption{数据集连续型特征的描述性统计（各特征可用样本量见 $n$ 列）}
  \label{tab:senp_feature_statistics}
  \begin{tabular}{lccccccc}
    \hline
    特征 & $n$ & 均值 & 标准差 & 中位数 & Q1 & Q3 & 偏度 \\
    \hline
    粒径 (nm)            & 339 & 65.30  & 79.72  & 30.00  & 30.00  & 92.15   & 3.92  \\
    水合粒径 (nm)        & 279 & 104.97 & 46.73  & 91.00  & 84.00  & 110.00  & 1.69  \\
    Zeta 电位 (mV)       & 144 & $-5.19$ & 29.76  & $-7.50$ & $-29.10$ & 29.10 & $-0.08$ \\
    暴露时间 (h)         & 363 & 55.21  & 26.14  & 72.00  & 24.00  & 72.00   & $-0.42$ \\
    浓度 ($\mu$g/mL)     & 363 & 12.59  & 38.17  & 0.90   & 0.52   & 10.00   & 8.28  \\
    细胞存活率 (\%)      & 371 & 49.43  & 28.50  & 50.00  & 41.00  & 68.00   & $-0.17$ \\
    \hline
  \end{tabular}
\end{table}

\begin{figure}[H]
  \centering
  \begin{minipage}{0.49\linewidth}
    \centering
    \begin{tikzpicture}
      \begin{axis}[
        acs ybar,
        width=\linewidth,
        height=5.8cm,
        ymin=0,
        ymax=405,
        ylabel={可用样本量 $n$},
        symbolic x coords={粒径,水合粒径,Zeta电位,暴露时间,浓度,细胞存活率},
        xtick=data,
        x tick label style={rotate=45, anchor=east, font=\scriptsize},
        title={\textbf{(A)} 数据覆盖度}
      ]
        \addplot+[fill=acsblue, draw=none] coordinates {(粒径,339) (水合粒径,279) (Zeta电位,144) (暴露时间,363) (浓度,363) (细胞存活率,371)};
      \end{axis}
    \end{tikzpicture}
  \end{minipage}
  \hfill
  \begin{minipage}{0.49\linewidth}
    \centering
    \begin{tikzpicture}
      \begin{axis}[
        acs axis,
        width=\linewidth,
        height=5.8cm,
        ymin=-1,
        ymax=9.2,
        ylabel={偏度},
        symbolic x coords={粒径,水合粒径,Zeta电位,暴露时间,浓度,细胞存活率},
        xtick=data,
        x tick label style={rotate=45, anchor=east, font=\scriptsize},
        title={\textbf{(B)} 分布偏斜}
      ]
        \addplot+[only marks, acsred, mark=square*, mark size=1.8pt] coordinates {(粒径,3.92) (水合粒径,1.69) (Zeta电位,-0.08) (暴露时间,-0.42) (浓度,8.28) (细胞存活率,-0.17)};
        \addplot+[black!55, dashed, mark=none] coordinates {(粒径,0) (细胞存活率,0)};
      \end{axis}
    \end{tikzpicture}
  \end{minipage}

  \vspace{0.9em}

  \begin{minipage}{0.62\linewidth}
    \centering
    \begin{tikzpicture}
      \begin{axis}[
        acs axis,
        width=\linewidth,
        height=5.2cm,
        xmin=110,
        xmax=390,
        ymin=-1,
        ymax=9.2,
        xlabel={可用样本量 $n$},
        ylabel={偏度},
        title={\textbf{(C)} 覆盖度--偏度诊断}
      ]
        \addplot+[
          only marks,
          mark=*,
          mark size=1.9pt,
          acsorange,
          nodes near coords,
          nodes near coords style={font=\tiny, anchor=west, xshift=1.5pt},
          point meta=explicit symbolic
        ] coordinates {
          (339,3.92) [粒径]
          (279,1.69) [水合粒径]
          (144,-0.08) [Zeta电位]
          (363,-0.42) [暴露时间]
          (363,8.28) [浓度]
          (371,-0.17) [细胞存活率]
        };
        \addplot+[black!55, dashed, mark=none] coordinates {(110,0) (390,0)};
      \end{axis}
    \end{tikzpicture}
  \end{minipage}
  \caption{主要连续型特征的可用样本量与偏度分布。}
  \label{fig:senp-continuous-summary}
\end{figure}

类别型特征的频率分布同样用于辅助判断模型可能学习到的主要数据模式（图~\ref{fig:senp-categorical-frequency}）。从目前所考察的类别型特征分布来看，数据集涉及 42 种细胞系，结肠癌细胞 HCT 116 出现频率最高（98 次，25.6\%），其次为肝癌细胞 HepG2（41 次，10.7\%）、黑色素瘤细胞 A375（23 次，6.0\%）和乳腺癌细胞 MCF-7（23 次，6.0\%），前 10 种细胞系合计覆盖 71.0\% 的样本。
组织来源以结肠（28.7\%）、肝脏（16.7\%）和脑组织（8.9\%）为主。

表面修饰类型共 43 种，其中聚糖修饰（PTR）占比最高（127 次，33.2\%），其后依次为聚糖蛋白复合物 PR（18 次，4.7\%）和 PAMAM+RGD+siRNA 复合修饰（15 次，3.9\%）。
合成方法以化学还原法为主（352 次，占 91.9\%），毒性检测方法中最常见的是 MTT 法（206 次，53.8\%）。

\begin{figure}[H]
  \centering
  \begin{minipage}{0.49\linewidth}
    \centering
    \begin{tikzpicture}
      \begin{axis}[
        acs xbar,
        width=\linewidth,
        height=5.8cm,
        xmax=215,
        xlabel={样本量},
        symbolic y coords={其他细胞系,MCF-7,A375,HepG2,HCT 116},
        ytick=data,
        y tick label style={font=\scriptsize},
        title={\textbf{(A)} 主要细胞系}
      ]
        \addplot+[fill=acsblue, draw=none] coordinates {(198,其他细胞系) (23,MCF-7) (23,A375) (41,HepG2) (98,HCT 116)};
      \end{axis}
    \end{tikzpicture}
  \end{minipage}
  \hfill
  \begin{minipage}{0.49\linewidth}
    \centering
    \begin{tikzpicture}
      \begin{axis}[
        acs axis,
        width=\linewidth,
        height=5.8cm,
        ymin=0,
        ymax=65,
        ylabel={样本占比（\%）},
        symbolic x coords={PTR,PR,PAMAM+RGD+siRNA,其他修饰},
        xtick=data,
        x tick label style={rotate=35, anchor=east, font=\tiny},
        nodes near coords,
        nodes near coords style={font=\tiny, yshift=1pt},
        title={\textbf{(B)} 表面修饰类型}
      ]
        \addplot+[only marks, acsorange, mark=square*, mark size=2pt] coordinates {(PTR,33.2) (PR,4.7) (PAMAM+RGD+siRNA,3.9) (其他修饰,58.2)};
      \end{axis}
    \end{tikzpicture}
  \end{minipage}
  \caption{主要细胞系与表面修饰类型的频率分布。}
  \label{fig:senp-categorical-frequency}
\end{figure}

因此，本研究的可视化分析并不只用于展示变量分布，而是与特征筛选和模型评价相结合：先通过统计图和表格检查样本覆盖范围、异常高剂量和类别集中度，再根据模型表现选择需要解释的关键特征，并在具备完整模型预测结果后利用 PDP、ICE 和双特征等高线图分析浓度、粒径、Zeta 电位、暴露时间、检测方法及材料化学增强特征等变量对细胞毒性预测的单独效应和交互效应。


\section{结果与讨论}

\paragraph{多类型纳米颗粒基准数据集上的模型比较。}
为获得可靠的性能基线，我们首先在大规模多类型纳米颗粒毒性数据集上比较六类模型的表现，包括随机森林\cite{breiman2001random}、XGBoost\cite{chen2016xgboost}、ABRF\cite{gromov2022attention}、LARF\cite{gromov2022larf}、TabTransformer\cite{huang2020tabtransformer} 以及蒸馏 TabTransformer\cite{hinton2015distilling}。该基准数据集来源于 Masarkar 等人建立的公开纳米颗粒毒性数据资源 \cite{masarkar2024}，覆盖金属、金属氧化物、聚合物、碳基与无机纳米颗粒，清洗后共包含 4102 个有效样本，按照 7:3 划分为训练集 3281 条与测试集 821 条。输入特征包括浓度、暴露时间、材料类别、粒径、Zeta 电位、细胞类型与检测方法等。

回归任务结果如表~\ref{tab:benchmark_regression} 所示。TabTransformer 取得最高的决定系数，$R^2$ 为 0.768，MAE 为 10.30\%，RMSE 为 14.97\%。XGBoost 与随机森林表现接近，分别达到 $R^2=0.752$ 与 $R^2=0.749$。引入注意力机制后的 ABRF 与 LARF 未超过传统树模型，蒸馏 TabTransformer 的性能则明显下降。

\begin{table}[htbp]
  \centering
  \caption{多类型纳米颗粒基准数据集上的回归性能比较}
  \label{tab:benchmark_regression}
  \begin{tabular}{lccc}
    \hline
    模型 & $R^2$ & MAE (\%) & RMSE (\%) \\
    \hline
    TabTransformer & 0.768 & 10.30 & 14.97 \\
    XGBoost & 0.752 & 11.13 & 15.56 \\
    随机森林 & 0.749 & 10.81 & 15.71 \\
    ABRF & 0.719 & 11.87 & 16.55 \\
    LARF & 0.716 & 11.94 & 16.63 \\
    蒸馏 TabTransformer & 0.506 & 16.85 & 22.31 \\
    \hline
  \end{tabular}
\end{table}

分类任务方面，树模型表现更稳健，其中 XGBoost 的 AUC 最高，为 0.945；随机森林的准确率与 F1 值最好，分别为 84.9\% 与 0.840。该基准比较表明，在大样本异质表格数据上，深度表格模型和树模型都具有较强竞争力，但树模型在分类任务中更稳定、训练成本更低。

\begin{table}[htbp]
  \centering
  \caption{多类型纳米颗粒基准数据集上的分类性能比较}
  \label{tab:benchmark_classification}
  \begin{tabular}{lcccc}
    \hline
    模型 & Accuracy & F1 & AUC \\
    \hline
    随机森林 & 0.849 & 0.840 & 0.943 \\
    XGBoost & 0.849 & 0.836 & 0.945 \\
    TabTransformer & 0.845 & 0.835 & 0.939 \\
    ABRF & 0.828 & 0.817 & 0.921 \\
    LARF & 0.819 & 0.809 & 0.914 \\
    蒸馏 TabTransformer & 0.722 & 0.707 & 0.801 \\
    \hline
  \end{tabular}
\end{table}

\paragraph{硒纳米颗粒专用数据集上的模型比较。}
在硒纳米颗粒专用数据集上，我们采用贴近实际应用的三折交叉验证方案，并以基础特征与增强特征两种设置分别训练随机森林、XGBoost 与 TabTransformer。基础特征包括浓度、暴露时间、粒径、Zeta 电位、尺寸比、表面电荷、细胞类型、修饰基团和检测方法等；增强特征在此基础上进一步引入 IC50、元素 Se 含量、长期稳定性、Se 3d 结合能 以及晶型编码。考虑到 IC50 与细胞毒性终点之间存在较强关联，以下结果主要用于比较不同增强信息的统计贡献，其中含 IC50 的结果需结合其解释边界加以理解。

结果如表~\ref{tab:senp_main_results} 所示。基础特征下，随机森林与 XGBoost 已具备较好的回归能力，分别达到 $R^2=0.701$ 与 $R^2=0.751$；TabTransformer 在该中小样本场景下明显欠拟合，$R^2$ 仅为 0.003。引入增强变量后，树模型的拟合能力进一步改善：随机森林的三折 CV MAE 从 11.04\% 降至 7.63\%，$R^2$ 从 0.701 提升至 0.915；XGBoost 的三折 CV MAE 从 9.87\% 降至 7.20\%，$R^2$ 进一步提升至 0.957，为所有方案中数值最高。相比之下，TabTransformer 即使加入增强特征后也仅达到 $R^2=0.057$，说明在当前样本规模下，树模型对硒纳米颗粒表格数据更具优势。

方法部分所述的多模态融合策略，即特征拼接、交叉注意力加权与残差连接，应用于全部增强特征配置。在当前数据集上，注意力机制相较简单拼接未带来明显提升，这可能与样本量有限有关；但该机制为引入可选 IC50 描述符及动态调整联合特征表示提供了明确路径。该设计将增强特征中的材料本征化学信息与通用理化描述符整合，对应于上述增强配置的性能改善。

需要强调的是，上述增强设置包含 IC50，因此其性能提升既反映了材料化学信息的补充作用，也可能部分受益于 IC50 与毒性终点之间的天然相关性，因而更适合作为辅助分析结果，而非直接视为真实部署场景下的最终预测性能。

\begin{table}[htbp]
  \centering
  \caption{硒纳米颗粒数据集上基础特征与增强特征的模型性能}
  \label{tab:senp_main_results}
  \begin{tabular}{lccc}
    \hline
    模型 & 特征集 & 3折 CV MAE (\%) & $R^2$ \\
    \hline
    随机森林 & 基础 & 11.04 & 0.701 \\
    随机森林 & 增强 & 7.63 & 0.915 \\
    XGBoost & 基础 & 9.87 & 0.751 \\
    XGBoost & 增强 & 7.20 & 0.957 \\
    TabTransformer & 基础 & 24.13 & 0.003 \\
    TabTransformer & 增强 & 25.50 & 0.057 \\
    \hline
  \end{tabular}
\end{table}

\paragraph{特征效应。}
增强特征模型的特征重要性分析表明，IC50、元素 Se 含量、长期稳定性和Se 3d 结合能 对硒纳米颗粒细胞毒性预测结果具有较大贡献。其中，IC50 直接反映剂量--效应关系，因此更应被视为毒理响应相关辅助变量；元素 Se 含量与颗粒纯度及潜在离子释放背景相关，长期稳定性表征颗粒在分散体系中的聚集与氧化趋势，Se 3d 结合能 则反映表面化学态差异。由此可见，在单一材料体系的毒性建模中，真正具有先验解释意义的增益主要来自能够刻画材料本征化学属性的变量，而不能仅依据含 IC50 模型的高分数判断模型已具备充分的真实预测能力。

\paragraph{关键特征对纳米硒材料选择的指导意义。}

\paragraph{组合特征分析。}
为了识别真正有效的增强信息，我们对五个候选增强变量进行了自动化筛选，候选变量包括 IC50、元素 Se 含量（wt\%）、长期稳定性（周）、Se 3d 结合能（eV）与晶型编码。穷举 31 种非空组合并以三折交叉验证 MAE 为目标后，表现最优的组合为
\[
\{\text{IC50},\ \text{元素 Se 含量},\ \text{长期稳定性},\ \text{Se 3d 结合能}\}.
\]
%% \[
%% \{\mathrm{IC50},\ \mathrm{元素\ Se\ 含量},\ \mathrm{长期稳定性},\ \mathrm{结合能\ Se3d}\}.
%% \]
该组合的三折 CV MAE 为 7.60\%，与使用全部增强特征时的 7.63\% 基本一致，说明晶型编码在当前数据集中的独立贡献有限，而 IC50、Se 含量、稳定性与结合能提供了主要统计增益。需要进一步指出的是，由于该组合包含 IC50，因此其结果更适合作为“含响应相关辅助变量时模型可达到的上界参考”，而不宜简单等同于仅依赖先验理化参数即可实现的预测效果。

\paragraph{模型结果对真实纳米硒材料设计与筛选的指导意义。}

\paragraph{回归与分类任务的最佳模型选择。}
综合上述结果，在硒纳米颗粒专用数据集上，\textbf{XGBoost 增强特征模型}在回归任务中取得了最高数值表现（$R^2=0.957$，CV MAE=7.20\%），高于随机森林和 TabTransformer。在分类任务中，XGBoost 和随机森林在基准数据集上表现相当，XGBoost 的 AUC 略高（0.945 vs 0.943）。但鉴于该回归最优结果对应的增强设置包含 IC50，本文更倾向于将其视为机理增强和辅助分析模型，而非直接作为真实应用场景中的唯一核心预测工具。相较之下，不依赖响应相关辅助变量的基础特征模型及其材料化学增强版本，更适合作为后续外部验证和部署研究的主要比较基线。

\paragraph{与现有工作的比较。}
与 Masarkar 等人 \cite{masarkar2024} 的通用纳米颗粒毒性预测模型相比，本研究具有以下改进：
\begin{itemize}
  \item \textbf{针对单一材料（纳米硒）} 构建了专用数据集，并引入了元素 Se 含量、长期稳定性和Se 3d 结合能 等材料化学增强信息，同时比较了含 IC50 与不含 IC50 的设置差异。这说明单一材料体系中的材料本征化学属性具有重要建模价值，而含 IC50 的高分数结果则更适合作为辅助参考，而非直接与仅依赖先验理化参数的通用模型进行等价比较。
  \item \textbf{系统比较了深度表格模型与树模型}，发现 TabTransformer 在小样本场景（383 条）下严重欠拟合，而树模型（尤其是 XGBoost）表现稳健。这为今后小样本纳米毒性数据集建模提供了参考。
  \item \textbf{通过穷举搜索确定了最优增强特征子集}，剔除了晶型编码等冗余特征，避免了过拟合，同时增强了模型的可解释性。
\end{itemize}
与近期其他基于注意力机制的树模型（如 ABRF、LARF）相比，本研究未发现注意力机制在小规模表格数据上带来显著提升，反而增加了计算开销。因此，对于纳米毒性这类中等规模、混合型特征的数据，传统的梯度提升树（XGBoost）仍是性价比最高的选择。

\paragraph{结论。}
本研究建立的 XGBoost 增强特征模型在当前数据集上表现出较强的拟合能力（$R^2=0.957$），并可同时提供分类决策（高/低存活）。但由于该结果对应的增强设置包含 IC50，故更适合作为辅助分析和机理增强参考。相对而言，元素 Se 含量、长期稳定性与Se 3d 结合能 等材料化学特征所体现的增益，更能支持单一纳米材料毒性建模中的特征工程选择，并为后续外部验证与模型部署提供比较基础。

上述多模态融合框架为整合异质特征模态提供了明确路径；后续将利用独立实验数据开展外部验证，进一步评价其泛化能力。

\section{结论}

本文围绕纳米硒（SeNPs）细胞毒性预测任务，构建了由基准数据集模型比较、纳米硒专用建模以及增强特征自动筛选组成的研究流程。结果表明，在大规模异质纳米颗粒表格数据上，深度表格模型与树模型均能够完成有效预测；而在样本规模有限、材料类型相对固定的纳米硒场景中，树模型表现出更高的稳定性和适用性。

纳米硒专用数据集包含从 43 篇发表于 2008--2022 年的同行评议文献整理出的 383 条有效记录。基础理化特征能够提供初步预测能力，而元素 Se 含量、长期稳定性和 Se 3d 结合能等材料化学增强变量可进一步改善模型拟合能力。本研究采用的多模态融合框架通过基础与增强描述符拼接、交叉注意力加权以及残差连接整合异质特征，在 IC50 等辅助信息可用时动态突出相关描述符。

在所比较的模型中，XGBoost 在增强特征配置下取得最高回归表现。但该配置包含与响应终点相关的辅助变量 IC50，而非仅包含预测前可获得的先验理化描述符，因此其结果更适合作为机理分析与验证的辅助基准，不宜作为前瞻预测的实际性能指标。随机森林则在精度、稳健性和可解释性之间保持了较好平衡，是后续平台部署与工程实现的更合适候选。

特征筛选表明，表现较优的增强变量组合无需包含全部候选描述符；仅使用 IC50、元素 Se 含量、长期稳定性和 Se 3d 结合能即可达到接近完整增强特征集的精度。这提示应区分两类增强变量：材料本征化学描述符，包括元素 Se 含量、稳定性和结合能，在单一材料体系中具有可迁移、可积累的特征工程价值；IC50 等响应相关辅助变量则在历史实验数据可用时提供诊断参考。因此，后续数据采集与实验设计应优先测量能够用于前瞻筛选的材料本征化学属性。

总体而言，本研究形成了由数据整理、模型比较、特征筛选和多模态融合组成的纳米硒毒性预测研究路径。\textbf{XGBoost 增强特征模型}提供了较高的诊断参考指标（$R^2=0.957$），而材料化学增强随机森林更适合作为后续部署的基线。目前正在表征的自合成纳米硒数据将用于后续外部验证，以评价 IC50 不可用的前瞻场景下模型的泛化能力。这些结果为计算辅助纳米毒性评价和实验优先级筛选提供了基础。


\section*{Acknowledgements}

Please use ``The authors thank \ldots'' rather than ``The authors would like to
thank \ldots''.

\clearpage
\section*{补充材料}
\begingroup
\setcounter{table}{0}
\renewcommand{\thetable}{S\arabic{table}}
\small
\setlength{\tabcolsep}{3.5pt}
\setlength{\LTleft}{0pt}
\setlength{\LTright}{0pt}
\renewcommand{\arraystretch}{1.12}
\begin{longtable}{r >{\raggedright\arraybackslash}p{0.21\linewidth} >{\raggedright\arraybackslash}p{0.66\linewidth}}
\caption{纳米硒专用数据集的来源文献（43 篇，2008--2022 年）。}\label{tab:source-publications}\\
\toprule
序号 & 来源文献 & 文献题名与 DOI \\
\midrule
\endfirsthead
\multicolumn{3}{l}{表 S1（续）} \\
\toprule
序号 & 来源文献 & 文献题名与 DOI \\
\midrule
\endhead
\midrule
\multicolumn{3}{r}{续下页} \\
\endfoot
\bottomrule
\endlastfoot
% 中英文共用同一组文献键；题名、作者、年份和 DOI 均由 references.bib 提供。
1 & \citeauthor{chen_selenium_2008} (\citeyear{chen_selenium_2008})~\cite{chen_selenium_2008} & \citetitle{chen_selenium_2008}\newline\citefield{chen_selenium_2008}{doi} \\
\addlinespace[4pt]
2 & \citeauthor{li_reversal_2011} (\citeyear{li_reversal_2011})~\cite{li_reversal_2011} & \citetitle{li_reversal_2011}\newline\citefield{li_reversal_2011}{doi} \\
\addlinespace[4pt]
3 & \citeauthor{zheng_sialic_2011} (\citeyear{zheng_sialic_2011})~\cite{zheng_sialic_2011} & \citetitle{zheng_sialic_2011}\newline\citefield{zheng_sialic_2011}{doi} \\
\addlinespace[4pt]
4 & \citeauthor{liu_selenium_2012} (\citeyear{liu_selenium_2012})~\cite{liu_selenium_2012} & \citetitle{liu_selenium_2012}\newline\citefield{liu_selenium_2012}{doi} \\
\addlinespace[4pt]
5 & \citeauthor{wu_surface_2012} (\citeyear{wu_surface_2012})~\cite{wu_surface_2012} & \citetitle{wu_surface_2012}\newline\citefield{wu_surface_2012}{doi} \\
\addlinespace[4pt]
6 & \citeauthor{chen_surface_2012} (\citeyear{chen_surface_2012})~\cite{chen_surface_2012} & \citetitle{chen_surface_2012}\newline\citefield{chen_surface_2012}{doi} \\
\addlinespace[4pt]
7 & \citeauthor{yu_positive_2012} (\citeyear{yu_positive_2012})~\cite{yu_positive_2012} & \citetitle{yu_positive_2012}\newline\citefield{yu_positive_2012}{doi} \\
\addlinespace[4pt]
8 & \citeauthor{chen_peg-nanolized_2012} (\citeyear{chen_peg-nanolized_2012})~\cite{chen_peg-nanolized_2012} & \citetitle{chen_peg-nanolized_2012}\newline\citefield{chen_peg-nanolized_2012}{doi} \\
\addlinespace[4pt]
9 & \citeauthor{pi_pathway_2013} (\citeyear{pi_pathway_2013})~\cite{pi_pathway_2013} & \citetitle{pi_pathway_2013}\newline\citefield{pi_pathway_2013}{doi} \\
\addlinespace[4pt]
10 & \citeauthor{wu_induction_2013} (\citeyear{wu_induction_2013})~\cite{wu_induction_2013} & \citetitle{wu_induction_2013}\newline\citefield{wu_induction_2013}{doi} \\
\addlinespace[4pt]
11 & \citeauthor{feng_differential_2014} (\citeyear{feng_differential_2014})~\cite{feng_differential_2014} & \citetitle{feng_differential_2014}\newline\citefield{feng_differential_2014}{doi} \\
\addlinespace[4pt]
12 & \citeauthor{huang_selective_2013} (\citeyear{huang_selective_2013})~\cite{huang_selective_2013} & \citetitle{huang_selective_2013}\newline\citefield{huang_selective_2013}{doi} \\
\addlinespace[4pt]
13 & \citeauthor{li_functionalized_2013} (\citeyear{li_functionalized_2013})~\cite{li_functionalized_2013} & \citetitle{li_functionalized_2013}\newline\citefield{li_functionalized_2013}{doi} \\
\addlinespace[4pt]
14 & \citeauthor{zhang_enhancement_2013} (\citeyear{zhang_enhancement_2013})~\cite{zhang_enhancement_2013} & \citetitle{zhang_enhancement_2013}\newline\citefield{zhang_enhancement_2013}{doi} \\
\addlinespace[4pt]
15 & \citeauthor{jiang_gracilaria_2014} (\citeyear{jiang_gracilaria_2014})~\cite{jiang_gracilaria_2014} & \citetitle{jiang_gracilaria_2014}\newline\citefield{jiang_gracilaria_2014}{doi} \\
\addlinespace[4pt]
16 & \citeauthor{yu_ph-responsive_2014} (\citeyear{yu_ph-responsive_2014})~\cite{yu_ph-responsive_2014} & \citetitle{yu_ph-responsive_2014}\newline\citefield{yu_ph-responsive_2014}{doi} \\
\addlinespace[4pt]
17 & \citeauthor{zeng_construction_2015} (\citeyear{zeng_construction_2015})~\cite{zeng_construction_2015} & \citetitle{zeng_construction_2015}\newline\citefield{zeng_construction_2015}{doi} \\
\addlinespace[4pt]
18 & \citeauthor{liu_sequentially_2016} (\citeyear{liu_sequentially_2016})~\cite{liu_sequentially_2016} & \citetitle{liu_sequentially_2016}\newline\citefield{liu_sequentially_2016}{doi} \\
\addlinespace[4pt]
19 & \citeauthor{liu_facile_2016} (\citeyear{liu_facile_2016})~\cite{liu_facile_2016} & \citetitle{liu_facile_2016}\newline\citefield{liu_facile_2016}{doi} \\
\addlinespace[4pt]
20 & \citeauthor{yu_x-ray-responsive_2016} (\citeyear{yu_x-ray-responsive_2016})~\cite{yu_x-ray-responsive_2016} & \citetitle{yu_x-ray-responsive_2016}\newline\citefield{yu_x-ray-responsive_2016}{doi} \\
\addlinespace[4pt]
21 & \citeauthor{fu_rgd_2016} (\citeyear{fu_rgd_2016})~\cite{fu_rgd_2016} & \citetitle{fu_rgd_2016}\newline\citefield{fu_rgd_2016}{doi} \\
\addlinespace[4pt]
22 & \citeauthor{fang_targeting_2017} (\citeyear{fang_targeting_2017})~\cite{fang_targeting_2017} & \citetitle{fang_targeting_2017}\newline\citefield{fang_targeting_2017}{doi} \\
\addlinespace[4pt]
23 & \citeauthor{chen_efficient_2018} (\citeyear{chen_efficient_2018})~\cite{chen_efficient_2018} & \citetitle{chen_efficient_2018}\newline\citefield{chen_efficient_2018}{doi} \\
\addlinespace[4pt]
24 & \citeauthor{fang_dualtargeted_2018} (\citeyear{fang_dualtargeted_2018})~\cite{fang_dualtargeted_2018} & \citetitle{fang_dualtargeted_2018}\newline\citefield{fang_dualtargeted_2018}{doi} \\
\addlinespace[4pt]
25 & \citeauthor{huang_autophagy_2018} (\citeyear{huang_autophagy_2018})~\cite{huang_autophagy_2018} & \citetitle{huang_autophagy_2018}\newline\citefield{huang_autophagy_2018}{doi} \\
\addlinespace[4pt]
26 & \citeauthor{liu_highly_2018} (\citeyear{liu_highly_2018})~\cite{liu_highly_2018} & \citetitle{liu_highly_2018}\newline\citefield{liu_highly_2018}{doi} \\
\addlinespace[4pt]
27 & \citeauthor{wang_regulating_2016} (\citeyear{wang_regulating_2016})~\cite{wang_regulating_2016} & \citetitle{wang_regulating_2016}\newline\citefield{wang_regulating_2016}{doi} \\
\addlinespace[4pt]
28 & \citeauthor{hu_selenium_2019} (\citeyear{hu_selenium_2019})~\cite{hu_selenium_2019} & \citetitle{hu_selenium_2019}\newline\citefield{hu_selenium_2019}{doi} \\
\addlinespace[4pt]
29 & \citeauthor{zeng_potentiation_2019} (\citeyear{zeng_potentiation_2019})~\cite{zeng_potentiation_2019} & \citetitle{zeng_potentiation_2019}\newline\citefield{zeng_potentiation_2019}{doi} \\
\addlinespace[4pt]
30 & \citeauthor{huang_dual-targeting_2019} (\citeyear{huang_dual-targeting_2019})~\cite{huang_dual-targeting_2019} & \citetitle{huang_dual-targeting_2019}\newline\citefield{huang_dual-targeting_2019}{doi} \\
\addlinespace[4pt]
31 & \citeauthor{huang_highly_2019} (\citeyear{huang_highly_2019})~\cite{huang_highly_2019} & \citetitle{huang_highly_2019}\newline\citefield{huang_highly_2019}{doi} \\
\addlinespace[4pt]
32 & \citeauthor{zhang_systematic_2019} (\citeyear{zhang_systematic_2019})~\cite{zhang_systematic_2019} & \citetitle{zhang_systematic_2019}\newline\citefield{zhang_systematic_2019}{doi} \\
\addlinespace[4pt]
33 & \citeauthor{zheng_thermosensitive_2019} (\citeyear{zheng_thermosensitive_2019})~\cite{zheng_thermosensitive_2019} & \citetitle{zheng_thermosensitive_2019}\newline\citefield{zheng_thermosensitive_2019}{doi} \\
\addlinespace[4pt]
34 & \citeauthor{huang_chiralitydriven_2020} (\citeyear{huang_chiralitydriven_2020})~\cite{huang_chiralitydriven_2020} & \citetitle{huang_chiralitydriven_2020}\newline\citefield{huang_chiralitydriven_2020}{doi} \\
\addlinespace[4pt]
35 & \citeauthor{yang_lentinan-functionalized_2020} (\citeyear{yang_lentinan-functionalized_2020})~\cite{yang_lentinan-functionalized_2020} & \citetitle{yang_lentinan-functionalized_2020}\newline\citefield{yang_lentinan-functionalized_2020}{doi} \\
\addlinespace[4pt]
36 & \citeauthor{liu_rapid_2021} (\citeyear{liu_rapid_2021})~\cite{liu_rapid_2021} & \citetitle{liu_rapid_2021}\newline\citefield{liu_rapid_2021}{doi} \\
\addlinespace[4pt]
37 & \citeauthor{an_functionalized_2021} (\citeyear{an_functionalized_2021})~\cite{an_functionalized_2021} & \citetitle{an_functionalized_2021}\newline\citefield{an_functionalized_2021}{doi} \\
\addlinespace[4pt]
38 & \citeauthor{song_translational_2021} (\citeyear{song_translational_2021})~\cite{song_translational_2021} & \citetitle{song_translational_2021}\newline\citefield{song_translational_2021}{doi} \\
\addlinespace[4pt]
39 & \citeauthor{yang_atherosclerotic_2022} (\citeyear{yang_atherosclerotic_2022})~\cite{yang_atherosclerotic_2022} & \citetitle{yang_atherosclerotic_2022}\newline\citefield{yang_atherosclerotic_2022}{doi} \\
\addlinespace[4pt]
40 & \citeauthor{yang_functionalized_2021} (\citeyear{yang_functionalized_2021})~\cite{yang_functionalized_2021} & \citetitle{yang_functionalized_2021}\newline\citefield{yang_functionalized_2021}{doi} \\
\addlinespace[4pt]
41 & \citeauthor{liang_translational_2022} (\citeyear{liang_translational_2022})~\cite{liang_translational_2022} & \citetitle{liang_translational_2022}\newline\citefield{liang_translational_2022}{doi} \\
\addlinespace[4pt]
42 & \citeauthor{liu_specific_2022} (\citeyear{liu_specific_2022})~\cite{liu_specific_2022} & \citetitle{liu_specific_2022}\newline\citefield{liu_specific_2022}{doi} \\
\addlinespace[4pt]
43 & \citeauthor{du_highly_2022} (\citeyear{du_highly_2022})~\cite{du_highly_2022} & \citetitle{du_highly_2022}\newline\citefield{du_highly_2022}{doi} \\
\addlinespace[4pt]

\end{longtable}
\noindent \textit{注：}逐文献记录数尚待根据原始 383 条数据核实。书目信息由中英文共享的参考文献库提供。

\endgroup
\clearpage

%%%%%%%%%%%%%%%%%%%%%%%%%%%%%%%%%%%%%%%%%%%%%%%%%%%%%%%%%%%%%%%%%%%%%
%% If you are using classical BibTeX rather than biblatex,
%% remove the \printbibliography[title={参考文献}] and uncomment the \bibliography one
%%%%%%%%%%%%%%%%%%%%%%%%%%%%%%%%%%%%%%%%%%%%%%%%%%%%%%%%%%%%%%%%%%%%%
\printbibliography
%\bibliography{acs-template.bib}

\newpage

\rule{0.05in}{1.75in}%
\begin{minipage}[b][1.75in]{3.25in}
  \sffamily
  \frenchspacing

  Some journals require a graphical entry for the Table of Contents. This
  should be laid out ``print ready'' so that the sizing of the text is correct.

  The space available depends on the journal: J. Am. Chem. Soc. allows 3.25 in
  by 1.75 in and requires sanserif text. Some journals want different sizes:
  you can easily adjust here.
  
  The two rules either side of the content are there to help judge the height
  of your material: they may be deleted once not required.
  
\end{minipage}%
\rule{0.05in}{1.75in}

\end{document}
